% page 490 du fac-simile (image p-489 du scan)
% bloc 162.6 x 250.6 mm ; marge gauche 8.0 mm
\pgc{-12.700mm}{0.000mm}{
\l{\cel{2}482}
\l{}
\l{\sou{rango}. I.Singla de la linei sur qui kozi, o personi, qui}
\l{\cel{3}intersucedas esas situita. - II. (a) En serio de personi,}
\l{\cel{3}de kozi, la plaso qua apartenas a singla de li, avan o}
\l{\cel{3}dop la ceteri. -b) La plaso quan ulu okupas en la esti-}
\l{\cel{3}mo da la cetera homi. - DEFR.}
\l{}
\l{\sou{rankorar}. (\sou{netrans., ad, pro, pri}).\cel{2}Memorar tenace la mala}
\l{\cel{3}agi quin onu subisis de ulu. - EFIS.}
\l{}
\l{\sou{ransono}. Preco postulata po la livro di milito-kaptito, di}
\l{\cel{3}karcerano. - DEF.}
\l{}
\l{\sou{ranunkulo}. (\sou{bot.}) Planto herbatra, di qua un speco, kun flo-}
\l{\cel{3}ri dazlante flava, esas vulgara en la prati. - L.\cel{1}\sou{ranun}-}
\l{\cel{3}\sou{culus.}\cel{2}- DEFIRS.}
\l{}
\l{\sou{ranuro}. (\sou{tekn.)}\cel{2}Kanelo facita en ligno-bloko por uzesar kom}
\l{\cel{3}kuliso o por recevar asemblo-peco. -\cel{2}FS.}
\l{}
\l{\sou{rapo}. (\sou{bot.}) Varietato di kaulo, kun radiko tre pulpoza e man-}
\l{\cel{3}jebla. - L.\cel{1}\sou{rapa}. - FIR.}
\l{}
\l{\sou{rapecar}.\cel{2}(\sou{trans.)}\cel{3}Igar bona-standa itere (stofo) per pasi-}
\l{\cel{3}gar interkruce fili en o sur la loko domajita (per truo,}
\l{\cel{3}per lacereso partala). - FI.}
\l{}
\l{\sou{rapida.}\cel{2}I. Qua parkuras multa spaco en poka tempo. - II.}
\l{\cel{3}Qua agas o facas en poka tempo (to pri quo onu esas paro-}
\l{\cel{3}lanta). - DEFIS.}
\l{}
\l{\sou{rapiero.}\cel{2}(\sou{skermarto}).Olim : Espado longa, akutigita, di qua}
\l{\cel{3}la konko esis truoza (trui aden qui povis penetrar la es-}
\l{\cel{3}pado di la adverso). - DEFR.}
\l{}
\l{raportar. (\sou{netrans., pri}). I. Naracar ad ulu to quon onu vidis,}
\l{\cel{3}audis, e c. - II. (\sou{metaf.})\cel{1}\sou{Raporto}\cel{1}: Rezultajo de la inter-}
\l{\cel{3}komparo di du quanti. - DEFIRS.}
\l{}
\l{\sou{rapsodo}. (en\cel{1}\sou{Grekia, antique)}.\cel{2}La recitisto di epikaji.- DEFIRS}
\l{}
\l{raptar. (\sou{trans., de)}.\cel{2}Deprenar per violento (ulu, o to quo}
\l{\cel{3}apartenas a kunhomo). - FIS.}
\l{}
\l{rara. Tre ne-ofta. DEFIS.}
\l{}
\l{\sou{raso.}\cel{2}I. Ensemblo ek la ancestri edek la decendinti di un}
\l{\cel{3}familio, di un populo. - II. Grupo di speco de animali}
\l{\cel{3}o de planti di qua la karakterizivo esas konstanta e trans-}
\l{\cel{3}misesas per genitado. - DEFIRS.}
\l{}
\l{raskalo. Persono qua nule skrupulas pri agar ne-honeste.- E.}
\l{}
\l{raspar. (\sou{trans.}) Transformar aden pulpo o pulvero, per ta im-}
\l{\cel{3}plemento koqueyala qua konsistas ye plako de lado herisita}
\l{\cel{3}ye asperaji sur qui onu frotas la substanco tale transfor-}
\l{\cel{3}menda.- DEFIRS.}
}
